\documentclass[10pt,a4paper]{amsart}
\usepackage[margin=19mm]{geometry}
\usepackage{amsmath,amssymb,mathtools}
\usepackage[scr=rsfs,cal=euler]{mathalfa}
\usepackage{xfrac}
\usepackage[unicode,hidelinks,bookmarks=false]{hyperref}
\newcommand{\ie}{i.e.\ }
\newcommand{\cf}{cf.\ }
\newcommand{\resp}{respectively\ }
\newcommand{\Subsection}{\subsection}
\providecommand{\nobreakdash}{\nobreak\hbox{-}}
\setlength{\emergencystretch}{2em}
\multlinegap0pt
\usepackage{bm}
\usepackage{bbm}
\usepackage{dsfont}
\usepackage{xspace}
\usepackage{cancel}
\usepackage{mathtools}
\usepackage{cleveref}
\usepackage{amsthm}
\makeatletter
\@ifundefined{theorem}{\newtheorem{theorem}{Theorem}}{}
\@ifundefined{cor}{\newtheorem{cor}[theorem]{Corollary}}{}
\@ifundefined{lemma}{\newtheorem{lemma}[theorem]{Lemma}}{}
\makeatother
\DeclareMathOperator{\pd}{pd}
\def\A{{\mathcal{A}}}% \A == \mathcal{A}
\def\B{{\mathcal{B}}}% \B == \mathcal{B}
\def\CC{{\mathbb{C}}}% \C == \mathbb{C}
\def\O{{\mathcal{O}}}% \G == \mathcal{G}
\title[Addition theorems for Ziegler pairs of hyperplane arrangemen]{Addition theorems for Ziegler pairs of hyperplane arrangements}
\author{Takuro Abe, Lukas K\"uhne, Piotr Pokora}
\date{Source: arXiv:2509.19011v2 (CC BY 4.0)}
\begin{document}
\maketitle

\noindent\emph{Source and licence.} Addition theorems for Ziegler pairs of hyperplane arrangements. Version arXiv:2509.19011v2 (CC BY 4.0), distributed under the Creative Commons Attribution 4.0 International licence, as recorded in the arXiv licence metadata for this exact version and shown on the abstract page https://arxiv.org/abs/2509.19011v2. Full rights notice: licenses/arxiv-2509.19011-CC-BY-4.0.txt.\par\smallskip

\noindent\textit{Editorial scope.} Counted unit: the Theorem whose source label is \texttt{thm:main\_general} in the article (source0.tex lines 592--603, source label \texttt{thm:main\_general}), delivered together with the article's own complete original proof of it (source0.tex lines 604--793, one \texttt{proof} environment of 190 lines). No mathematics is written, repaired or re-derived: statement and proof are the article's own text line for line. Required context retained uncounted: lem:surjectivity (lines 360--368); cor:generic (lines 427--437); main (lines 459--463). Every definition, display and label of those lines is retained unchanged. Omitted from this package and not redistributed: 1--359, 369--426, 438--458, 464--591, 794--859. Nothing inside a retained excerpt is omitted or altered: every retained body is delivered exactly as the article's lines are written, so no omission receipt of any kind accompanies this package. Citations: the retained excerpts carry no citation command at all, so no bibliography entry of the article is transcribed and no reference is redistributed. Reference adaptation: none was needed and none was made; every reference inside the delivered unit resolves inside it, so no unresolved reference or citation is produced. Printed slips kept literal: no printed word, symbol, hypothesis or number of the counted statement or of its proof was altered. Layout and class: the article's own class is \texttt{amsart} with packages including latexsym, xy, amssymb, amsmath, amsthm, thmtools, biblatex, hyperref, cleveref, tikz; the wrapper is the lane's pinned \texttt{amsart} editorial wrapper at 10pt on A4, which restates the theorem-like environments the retained text needs and adds the article's own macro definitions that the retained text needs (\texttt{\textbackslash A}, \texttt{\textbackslash B}, \texttt{\textbackslash CC}, \texttt{\textbackslash O}, \texttt{\textbackslash pd}), with extra wrapper packages (bm, bbm, dsfont, xspace, cancel, mathtools, cleveref). The delivered numbering is the unit's own. Assets: no retained span contains a figure environment, a graphics-inclusion command, a PDF-page inclusion, a table or a TikZ picture; no image file is copied into this package, the package declares no asset dependency and no image file is redistributed with it. Licence: version arXiv:2509.19011v2 of this article is distributed under the Creative Commons Attribution 4.0 International licence, as recorded in the arXiv licence metadata for this exact version and shown on the abstract page https://arxiv.org/abs/2509.19011v2. A full rights notice is delivered at licenses/arxiv-2509.19011-CC-BY-4.0.txt. Characters: every retained excerpt is pure ASCII, so the delivered engine, XeLaTeX with its default Latin Modern text font, reports no missing character.
\begin{lemma}\label{lem:surjectivity}
	Let $\A$ be an irreducible arrangement in $\CC^3$. Let $\A_\ell$ be the arrangement obtained from $\A$ by taking $(\ell-3)$-times the coning operation, i.e., $\A_\ell \subset \CC^\ell$ with new coordinates $x_4,\ldots,x_\ell$.
	Let $H\subset \CC^\ell$ be a generic hyperplane with respect to $\A_\ell$.
	Then for $\B:=\A_\ell \cup \{H\}$ the Euler restriction map
	\[
	\rho \widetilde{D(\B)} \twoheadrightarrow \widetilde{D(\B^H)}
	\]
	is a surjective map on the level of sheaves.
\end{lemma}

\begin{cor}\label{cor:generic}
	Let $\mathcal{A}$ be an arrangement of hyperplanes in $\CC^3$.
	Suppose that $\theta_E,\theta_2,\ldots,\theta_s$ are minimal generators of $D(\A)$.
	For $k=3,\dots,\ell$, let $\A_k$ be the arrangement obtained from $\A$ by taking $(k-3)$-times the coning operation.
	Let $H$ be a generic hyperplane with respect to $\A_\ell$ in $\mathbb{C}^\ell$ and set $H_k:= H \cap \bigcap_{j=k+1}^{\ell} \{x_j=0\}$ for $k=3,\dots,\ell$.
	Then it holds that
		\begin{enumerate}
			\item[(i)] $H_k$ is generic with respect to $\A_k$ for all $k=3,\dots,\ell$.
			\item[(ii)] $H_k$ is algebraically generic with respect to $\A_k$ and its minimal set of generators $\theta_E,\theta_2,\dots,\theta_s,x_4\partial_{x_4},\dots,x_k \partial_{x_k}$ for all $k=3,\dots,\ell$.
		\end{enumerate}
\end{cor}

\begin{theorem}\label{main}
Let \(\mathcal{A}_1,\mathcal{A}_2\) be central arrangements in \(\mathbb C^3\) with
\(L(\mathcal{A}_1)\cong L(\mathcal{A}_2)\) and \(\exp(\mathcal{A}_1)\ne \exp(\mathcal{A}_2)\).
Let \(\mathcal{B}_i=\mathcal{A}_i\cup\{H\}\), where \(H\) is a generic hyperplane, and assume \(L(B_1)\cong L(B_2)\). Then \(B_1,B_2\) form a Ziegler pair.
\end{theorem}

\begin{theorem}\label{thm:main_general}
	Let $\A$ be an irreducible arrangement in $\CC^3$. Let $\A_\ell$ be the arrangement obtained from $\A$ by taking $(\ell-3)$-times the coning operation, i.e., $\A_\ell \subset \CC^\ell$ with new coordinates $x_4,\ldots,x_\ell$.
	Let $H\subset \CC^\ell$ be a generic hyperplane with respect to $\A_\ell$.
	Then for $\B_\ell:=\A_\ell \cup \{H\}$ and $n:=|\A^{H_3}|-1$ we have
	$$
	\exp_0(\B_\ell)=((2)^{\binom{\ell-2}{2}},(\exp_0(\A)+ 1)^{\ell-2},n).
	$$

    %PP -  not sure about the number of generators, I guess we have (\ell-3)+\binom{\ell-3}{2}} = \binom{\ell-2}{2} generators in degree 2.
%	 for $k\ge 3$ we set $H_k:= H \cap \bigcap_{j=k+1}^{\ell} \{x_j=0\}$, which is a hyperplane in $\CC^k$, and $n+1:=|\A^{H_3}|$.
	\label{mainhigh}
\end{theorem}

\begin{proof}
Let $\theta_E,\theta_2,\ldots,\theta_s$ be a minimal set of generators for $D(\A)$ with $\deg \theta_i=d_i$. Then it is clear that $D(\A_\ell)$ has a minimal set of generators 
$$
\theta_E,\theta_2,\ldots,\theta_s,x_4\partial_{x_4},\ldots,x_\ell\partial_{x_\ell}.
$$

For $k\ge 3$, we define $\A_k$ to be the $(k-3)$--times coned arrangement of~$\A$ and
$\B_k:=\A_k \cup \{H_k\}$.
By the construction of the coned arrangement and the generic hyperplane $H$, we have
\[
	\B_\ell^{H_k}=\B_k^{H_k} = \B_{k-1}.
\]
Let $\rho_k^{k-1}:D(\B_k)\to D(\B_k^{H_k})=D(\B_{k-1})$ be the Euler restriction map, and for $2\le j < i \le \ell$ we denote by $\rho_i^j:D(\B_i)\to D(\B_j)$ the composition of these maps.

\emph{Claim 1:} For every $3\le k \le \ell$, the map $\rho_k^{k-1}:D(\B_k)\to D(\B_{k-1})$ is surjective and therefore the following Euler sequence is exact:
\begin{equation}\label{eq:exact_euler}
0 \rightarrow D(\A_k)[-1] \stackrel{\cdot \alpha_{H_k}}{\rightarrow}
D(\B_k) \stackrel{\rho_k^{k-1}}{\rightarrow } D(\B_{k-1}) \rightarrow 0.
\end{equation}
\emph{Proof of Claim 1:}
The case $k=3$ was proved in~\Cref{main}.
Now assume $k>3$.
Since $\pd_S D(\A_k)\le 1$, this module has a free resolution
$$
0 \rightarrow 
\bigoplus_{j=1}^{p - k} S[-b_j]
\rightarrow 
\bigoplus_{j=1}^{p} S[-a_j]
\rightarrow D(\A_k) \rightarrow 0
$$
for some numbers $a_j, b_j$ and $p=s+k-3$.
Since we assume $k>3$, we know that $H^1(\O(-a_j))=0$ for all $1\le j \le p$ and $H^2(\O(-b_j))=0$ for all $1\le j \le p-k$.
Thus the long exact sequence in cohomology corresponding to the above free resolution shows that $H^1(\widetilde{D(\A_k)(-1)})=0$.

Lemma~\ref{lem:surjectivity} implies that there is an exact sequence of sheaves
\[
0 \rightarrow 
\widetilde{D(\A_k)}(-1) \stackrel{\cdot \alpha_{H_k}}{\rightarrow}
\widetilde{D(\B_k)} \stackrel{\rho_k^{k-1}}{\rightarrow}
\widetilde{D(\B_{k-1})} \rightarrow 0.
\]
The above cohomology-vanishing argument now shows that taking global sections preserves exactness and thus the map $\rho_k^{k-1}$ is surjective, which proves Claim 1.

The module $D(\B_\ell^{H_3})$ is a free module of rank $2$ with a basis $\theta_E,\phi_2$ with $\deg \phi_2 = n$.
Since the map $\rho_k^{k-1}$ is surjective for every $k\ge 3$ we can find successive preimages of $\phi_2$, namely $\phi_k\in D(\B_k)$ such that $\rho_k^2(\phi_k)=\phi_2$ for every $k\ge 3$.

After a change of coordinates, we may assume that $\partial_{x_j}(\alpha_H)=1$ for all $4\le j \le \ell$.
For $2\le i \le s$ and $4\le j\le \ell$, consider the derivation
$$
\varphi_i^j:=x_j \theta_i-\theta_i(\alpha_H) (x_j \partial_{x_{j}}).
$$
Then $\varphi_i^j(\alpha_H)=0$ for all $ 2 \le i \le s,\ 
4 \le j \le \ell$, and $\varphi_i^j \in D(\B_\ell)$ with $\deg \varphi_i^j=d_i+1$.
We also define the derivation for $4 \le i<j$
$$
\eta_i^j:=x_ix_j(\partial_{x_i}-\partial_{x_j}).
$$
Then clearly $\eta_i^j \in D(\B_k)$ for $j \le k$.

For $k\ge 3$, consider the set of derivations
\begin{multline*}
	G_k := \bigg\{\theta_E,\{\alpha_{H_k} x_j \partial_{x_j}\}_{4\le j \le k}, \{\alpha_{H_k} \theta_i\}_{i=2}^s, 
	\{ \varphi_i^j\}_{2 \le i \le s,\ 4\le j \le k}, \\ \{\eta_i^j\}_{4 \le i<j\le k}, \phi_k \bigg\}.
\end{multline*}

\emph{Claim 2:} For each $k\ge 3$, the set $G_k$ generates $D(\B_k)$.

\emph{Proof of Claim 2:} We proceed by induction, where the base case $k=3$ was again proved in~\Cref{main}.
Now assume the claim holds for some $3\le k-1$.

The derivations $\theta_i$ only involve the variables $x_1,x_2,x_3$, and by Corollary~\ref{cor:generic}, it holds that $\theta_i(\alpha_{H_k})\notin S\alpha_{H_k}$ for all $2 \le i \le s$ and $3\le k \le \ell$.
Therefore $\rho_k^{k-1}(\varphi_{i}^j)=\varphi_{i}^j \in D(\B_{k-1})$ for all $4< k, j \le \ell$ and $2 \le i \le s$. Now let us check the images of $\varphi_i^k$ and $\eta_i^k$ under the map $\rho:=\rho_k^{k-1}$.

First, we show that $\rho(\varphi_i^j)=\varphi_i^j$ if $j<k$, and 
$\rho(\varphi_i^k)=-\alpha_{H_{k-1}}\theta_i$.
For $j<k$, there are no $x_k$-terms in $\varphi_i^j$. Thus 
$\rho(\varphi_i^j)=\varphi_i^j$. Let $j=k$, now for $t<k$ 
\begin{eqnarray*}
\rho(\varphi_i^k)(x_t)&=&
x_k \overline{\theta_i(x_t)}\\
&=&-\alpha_{H_{k-1}}\overline{\theta_i(x_t)}=(-\alpha_{H_{k-1}} \theta_i)(x_t)
\end{eqnarray*}
since modulo $x_k$ does not change $\theta_i$. Moreover, 
\begin{eqnarray*}
\rho(\varphi_i^k)(x_k)&=&
-\overline{\theta_i(\alpha_H)} x_k\\
&=&-\alpha_{H_{k-1}}\theta_i(-\alpha_{H_{k-1}}).
\end{eqnarray*}
Since $\overline{x_k}=-\alpha_{H_{k-1}}$, we know that 
$\rho(\varphi_i^k)=-\alpha_{H_{k-1}}\theta_i$.

Next, let us show that $\rho(\eta_i^k)=-\alpha_{H_{k-1}} x_i\partial_{x_i}$.
Clearly $\rho(\eta_i^k)(x_t)=0$ if $t \neq i,k$. We compute 
\begin{eqnarray*}
\rho(\eta_i^k)(x_i)&=&=\overline{x_ix_k}\\
&=&-\alpha_{H_{k-1}} x_i=
(-\alpha_{H_{k-1}} x_i \partial_{x_i})(x_i).
\end{eqnarray*}
Furthermore,
\begin{eqnarray*}
\rho(\eta_i^k)(x_k)&=&\overline{x_i\alpha_{H_{k-1}}}\\
&=&\alpha_{H_{k-1}} x_i=
(-\alpha_{H_{k-1}} x_i \partial_{x_i})(-\alpha_{H_{k-1}}),
\end{eqnarray*}
which means that $\rho(\eta_i^k)=-\alpha_{H_{k-1}} x_i\partial_{x_i}$.
We have also $\rho(\eta_i^j)=\eta_i^j$ for  $4\le i < j <k$.
 Hence, by $\rho_k^{k-1}$, 
$$
\{\varphi_i^j\}_{2 \le i \le s,4 \le j \le k} \mapsto
\{\varphi_i^j,\alpha_{H_{k-1}} \theta_i\}_{2 \le i \le s,4 \le j \le k-1}
$$
and 
$$
\{\eta_i^k\}_{4 \le i \le k-1} \mapsto 
\{\alpha_{H_{k-1}} x_i \partial_{x_i}\}_{4 \le i \le k-1}
$$
up to signs, so we arrive at $\rho_k^{k-1}(G_k)=G_{k-1}$.
Using the exact sequence~\eqref{eq:exact_euler}, we obtain a generating set of $D(\B_k)$ by mapping the generators of $D(\A_k)$ and selecting preimages of the generators of $D(\B_{k-1})$ under the surjective map $\rho_k^{k-1}$.
The latter module is generated by $G_{k-1}$ by our induction assumption.

\emph{Claim 3:} The set $G_k$ is a minimal generating set of $D(\B_k)$. Since the degrees of the set $G_k$ for $k=\ell$ match the claimed exponents, this completes the proof the theorem once this claim is proved.

\emph{Proof of Claim 3:} We again proceed by induction where the base case $k=3$ was proved in~\Cref{main}.
Suppose that $G_{k-1}$ minimally generates $\B_{k-1}$ for some $3\le k-1$.

Since $\rho_k^{k-1}$ is surjective, we cannot remove 
$\theta_E$, $\varphi_i^j$ for $2 \le i \le s, 4 \le j \le k$, $\eta_i^j$ for $4\le i < j \le k$ or $\phi_k$.
Also, since $\A$ is irreducible, the derivations $\theta_2,\dots,\theta_s$ are all of degree at least $2$.
Thus the derivations $\varphi_i^j,\alpha_H \theta_i$  have degree at least $3$.
As $\A$ contains at least $4$ hyperplanes, also the derivation $\phi_k$ has degree at least $3$.

Suppose $x_t\alpha_{H_k} \partial_{x_t}$ is removable, then we have a relation involving the other derivations in $G_k$ of degree $2$, namely
\[
	x_t\alpha_H\partial_{x_t}=
	\sum_{i \neq t} c_ix_i\alpha_{H_k}\partial_{x_i}
	+\sum_{4\le i < j \le k} e_{ij} x_i x_j(\partial_{x_i}-\partial_{x_j}),
\]
where $c_i$ and $e_{ij}$ are constants.
Restricting this equation to $x_k=0$ yields a linear relation among the generators $G_{k-1}$, which is a contradiction to our induction assumption.

Moreover, suppose we can express $\alpha_{H_k}\theta_t$, for $1\le t \le s$, via the relation
\begin{equation}\label{eq:relation}
	\alpha_{H_k}\theta_t=f \theta_E+\sum_{i\neq t} f_i \alpha_{H_k} \theta_i+
	\sum_{\substack{1\le i \le s, \\4\le j \le k}}g_{ij} \varphi_i^j+\sum_{i=4}^k h_{i} x_i\alpha_{H_k} \partial_{x_i}+
	\sum_{4\le i < j \le k} q_{ij}\eta_i^j+
	c \phi_k,
\end{equation}
for polynomials $f, f_i, g_{ij},h_i,q_{ij}\in \mathbb{C}[x_1,\ldots,x_k]$.

The regularity argument from above tells us that for each $i=2,\ldots,s$ the degree of $\theta_i$ is at most $|\A|-2$.
As $\phi_k$ has degree $|\A|-1$, we can conclude that $c$ must be a constant.
Restricting this relation to $H_3$, i.e., $\alpha_{H_k}=\alpha_{H_{k-1}}=\dots = \alpha_{H_{3}}$, yields the relation
\[
0 = \overline{f} \theta_E + \overline{c} \phi_2.
\]
Since these two derivations form a basis of $D(\B_2)$ this yields $c=0$.

%Note that $\varphi_i^j|_{x_\ell=0}=0$ by construction, $\theta_i|_{x_\ell=0}=\theta_i$ and $\alpha_H|_{x_\ell=0}=\alpha_{H_{\ell-1}}$. So restricting the above onto $x_\ell=0$, we have 

%So the rest is to show that they are minimal. 
%Since $\rho_\ell^{\ell-1}$ is surjective, we cannot remove 
%$\theta_E,\varphi_i^j,\eta_i^\ell,\phi$.
%Also, since $\B_\ell$ is irreducible, $\deg \varphi_i^j,\alpha_H \theta_i,\phi$ has degrees at least three. So if $x_i\alpha_H \partial_{x_i}$ is removable, then say $i=t,$
%$$
%x_t\alpha_H\partial_{x_t}=
%\sum_{i \neq t} c_ix_i\alpha_H\partial_{x_i}
%+\sum_{i=4}^{\ell-1} e_i x_\ell x_i(\partial_{x_i}-\partial_{x_\ell}).
%$$
%Here $c_i,e_i$ are all constants. However, by putting $x_\ell=0$, we have a contradiction. So we cannot remove $x_i\alpha_H \partial_{x_i}$.
%Finally, assume that 
%$$
%\alpha_H \theta_1=f \theta_E+\sum_{i=2}^s f_i \alpha_H \theta_i+
%\sum_{i,j}g_{ij} \varphi_i^j+\sum_{i=4}^\ell h_{i} x_i\alpha_H \partial_{x_i}+
%\sum q_{i}\eta_i^\ell+
%c \phi.
%$$
%By constructions, it is clear that, if we restrict this set of generators onto $H_i$, then we obtain the set of minimal generators for $D(\A_i)$ by induction. 
%We may assume that $f=0$ by replacing, say $\phi$ by $\phi+f\theta_E$.
%Assume that $c \neq 0$. By the regularity argument, we know that $\deg \theta_i,\deg \varphi_i^j$ is at most $|\A_3|-1$. Since $\deg \phi=|\A_3|-1$, it holds that $c$ is a constant. Then by restricting the equation onto $H_3$, we know that $0=c\phi$, a contradiction. So $c=0$. 
Note that, by construction, $\varphi_i^k|_{x_k=0}=0$  and $\varphi_i^j|_{x_k=0}=\varphi_i^j$ for $j<k$.
The same conclusion holds for the restriction of $\eta_i^j$.
 Furthermore, it holds that $\theta_i|_{x_k=0}=\theta_i$ and $\alpha_{H_k}|_{x_k=0}=\alpha_{H_{k-1}}$.
 Now restricting the relation~\eqref{eq:relation} onto $x_k=0$ yields
 \begin{align*}
 		\alpha_{H_{k-1}}\theta_t=&\overline{f} \theta_E+\sum_{i\neq t} \overline{f_i} \alpha_{H_{k-1}} \theta_i+\\
 	&\sum_{\substack{1\le i \le s, \\4\le j \le k-1}}g_{ij} \varphi_i^j+\sum_{i=4}^{k-1} \overline{h_{i}} x_i\alpha_{H_{k-1}} \partial_{x_i}+
 	\sum_{4\le i < j \le k-1} \overline{q_{ij}}\eta_i^j.
 \end{align*}
 This is impossible as these are all derivations in $G_{k-1}$ and hence form a minimal generating set of $D(\B_{k-1})$ by our induction hypothesis.
\end{proof}

\end{document}
